\subsection{静态区间求和（教学例题）}

\practice{P8218}{【深进1.例1】求区间和}
使用一维前缀和；每次查询计算右端前缀减去左端前一项。

输入长度为 $n$ 的数组和 $q$ 次闭区间查询，每次输出区间和。$1\le n,q\le2\times10^5$，$|a_i|\le10^9$，保证 $1\le l\le r\le n$。第一行 $n,q$，第二行数组，随后 $q$ 行各给出 $l,r$。
\begin{verbatim}
输入
5 3
2 -1 3 4 5
1 3
2 4
5 5
输出
4
6
5
\end{verbatim}
前缀数组为 $[0,2,1,4,8,13]$，查询 $[2,4]$ 得 $8-2=6$。0号位置让 $l=1$ 的查询无需特殊判断。使用 long long 存前缀和。总复杂度 $O(n+q)$，空间 $O(n)$。
\Code{05_prefix.cpp}
\subsection{批量区间加法（教学例题）}

\practice{P2367}{语文成绩}
差分恢复全部成绩后，原题还要取最小值，而不是输出整个数组。

输入长度为 $n$ 的数组，执行 $q$ 次操作，每次给闭区间 $[l,r]$ 加 v，最后输出整个数组。$1\le n,q\le2\times10^5$，初始值和修改值绝对值不超过 $10^9$。
\begin{verbatim}
输入
5 2
1 2 3 4 5
2 4 3
1 5 -1
输出
0 4 5 6 4
\end{verbatim}
先用初始数组构造差分。所有修改都可叠加，最后做一次前缀累加。总时间 $O(n+q)$，空间 $O(n)$。上例第一轮得到 $[1,5,6,7,5]$，第二轮整体减1。
\Code{06_difference.cpp}
\subsection{切木条的最大整数长度（教学例题）}

\practice{P2440}{木材加工}
用“能切出的段数不少于所需数量”作判定，找最后一个可行长度。

有 n 根木条，第 i 根长度为 $a_i$。每根可切成若干长度相同的整段，剩余部分丢弃，不能拼接不同木条。至少要得到 k 段，求最大正整数段长；若连长度1都无法满足，输出0。$1\le n\le2\times10^5$，$1\le a_i\le10^9$，$1\le k\le10^{12}$。输入第一行 n、k，第二行各长度。
\begin{verbatim}
输入
4 7
20 15 10 17
输出
7
\end{verbatim}
长度7时可得 $2+2+1+2=7$ 段；长度8时只有 $2+1+1+2=6$ 段。判定函数为 $\sum_i\lfloor a_i/x\rfloor\ge k$，随 x 增大计数不增。0作为“无正整数解”的哨兵，不传入判定函数，因此不会除以0。复杂度 $O(n\log\max a_i)$，空间 $O(n)$。
\Code{07_wood.cpp}
\subsection{乘积小于k的连续区间个数（教学例题）}

给定 n 个\textbf{正整数}，统计乘积严格小于 k 的非空连续子数组数量。$1\le n\le2\times10^5$，$1\le a_i\le10^6$，$0\le k\le10^9$。输入第一行 n、k，第二行数组。
\begin{verbatim}
输入
4 100
10 5 2 6
输出
8
\end{verbatim}
合法区间包括4个单元素、$[10,5]$、$[5,2]$、$[2,6]$、$[5,2,6]$，共8个。维护窗口乘积，每加入右端点，就在乘积过大时删除左端点。缩小后，以当前右端点结尾的合法区间有 $right-left+1$ 个。

所有数至少为1，因此缩短区间不会使乘积增大，合法左端点形成后缀。若允许负数，这个性质会失效；若允许0，除法维护也需另外处理。k不大于1时没有合法非空区间。按本题约束，乘入新元素前乘积小于 $10^9$，乘后小于 $10^{15}$，long long 足够。时间 $O(n)$，空间 $O(n)$。
\Code{08_window.cpp}
\subsection{合并石堆的最小代价（教学例题）}

有 n 堆石子，每次可以任意选两堆合并，代价等于两堆石子数之和。求合成一堆的最小总代价。本题\textbf{没有相邻限制}。$1\le n\le2\times10^5$，每堆石子数在1到 $10^9$ 之间。输入第一行 n，第二行各堆大小。
\begin{verbatim}
输入
4
1 2 3 4
输出
19
\end{verbatim}
先合并1和2，代价3；再合并3和3，代价6；最后合并4和6，代价10。总计19。

\textbf{为什么取最小两堆。}合并过程可以看成二叉树，每个原始堆是叶子；总代价等于每个叶子权值乘以其深度之和。在固定树形中，较小权值放到较深位置不会使代价增大。选择一对最深的兄弟叶子，可通过交换把最小的两个权值放到那里。因此存在一个最优方案先合并最小两堆；把它们缩成一堆后，剩下同类子问题。重复此选择即可。

用小根堆维护所有堆，每次取两个最小值再插入它们的和。时间 $O(n\log n)$，空间 $O(n)$。只有一堆时无需合并，答案为0。
\Code{09_merge.cpp}
\subsection{按字典序输出全排列（教学例题）}

\practice{P1706}{全排列问题}
原题要求每个数占5个字符宽度，应使用\texttt{printf("\%5d", x)}，不能直接沿用单空格输出。

输入 $1\le n\le8$，输出1到n的所有排列，每行一个，按字典序升序。
\begin{verbatim}
输入
3
输出
1 2 3
1 3 2
2 1 3
2 3 1
3 1 2
3 2 1
\end{verbatim}
path[0..depth-1] 保存已选的前缀，used[x] 标记 x 是否已在前缀中。每层从小到大尝试尚未使用的数，输出顺序自然为字典序。每个排列对应唯一一条长度 n 的路径，因此不重不漏。输出本身有 $n\cdot n!$ 个数字量级，时间 $O(n\cdot n!)$，递归与标记空间 $O(n)$。
\Code{10_permutations.cpp}
\subsection{迷宫最短步数（教学例题）}

\practice{P1443}{马的遍历}
将四方向改成马的八种跳法，按原题输出整个距离矩阵；不是只输出终点距离。

输入 $n\times m$ 字符网格，点号为可走格，井号为障碍。每步上下左右移动一格，求左上角到右下角的最少步数；不可达输出-1。$1\le n,m\le500$，起终点也可能是障碍。第一行 n、m，随后 n 行长度为 m 的字符串。
\begin{verbatim}
输入
3 4
....
.##.
....
输出
5
\end{verbatim}
\textbf{解答。}所有移动代价相等，使用 BFS。起点距离为0，邻居距离为当前距离加1；入队时就记录距离，确保同一格最多入队一次。队列按距离层次扩展，因此首次到达就是最短距离。注意这里的“步数”不包含起点，单格可走迷宫答案为0。时间与空间均为 $O(nm)$。
\Code{11_bfs.cpp}
\subsection{最长严格上升子序列（教学例题）}

\practice{B3637}{最长上升子序列}
先掌握以当前位置结尾的状态，再根据数据范围选择二次或对数优化。

输入 n 个整数，选择一些位置并保持原有顺序，使选出的值严格递增，求最大长度。$1\le n\le5000$，值的绝对值不超过 $10^9$。子序列可以不连续。第一行 n，第二行数组。
\begin{verbatim}
输入
6
3 1 2 2 5 4
输出
3
\end{verbatim}
可取 $[1,2,5]$ 或 $[1,2,4]$。设 f[i] 为\textbf{以第i个元素结尾}的最长长度，
\[f[i]=1+\max\{f[j]:j<i,\ a_j<a_i\},\]
没有合法 j 时 f[i]=1。任何以 i 结尾的较长方案，倒数第二个位置必然是一个合法 j；反过来，把 a[i] 接到该 j 的方案后面也合法。因此枚举所有前驱就覆盖全部可能。答案是所有 f[i] 的最大值，不一定是 f[n]。时间 $O(n^2)$，空间 $O(n)$。
\Code{12_lis.cpp}
\subsection{非空最大子段和（教学例题）}

\practice{P1115}{最大子段和}
必须选择非空连续区间，检查全负数组。

输入 n 个整数，求一个\textbf{连续且非空}区间的最大和。$1\le n\le2\times10^5$，$|a_i|\le10^9$。
\begin{verbatim}
输入
5
-2 3 -1 4 -5
输出
6
\end{verbatim}
设 ending 为以当前位置结尾的非空区间最大和。若使用当前数 x，则最优区间要么从 x 重新开始，要么延长前一位置的最优区间：$ending=\max(x,ending+x)$。全局答案记录所有位置的最大 ending。由于要求非空，初始化为第一个元素，不能初始化为0。时间 $O(n)$，额外空间 $O(1)$。
\Code{13_maxsum.cpp}
\subsection{数字三角形的最大路径和（教学例题）}

输入 n 行数字，第 i 行有 i 个数。从顶点出发，每次走到下一行的正下方或右下方，求到最后一行的最大路径和。$1\le n\le1000$，每个数的绝对值不超过 $10^6$。第一行 n，随后各行给出三角形。
\begin{verbatim}
输入
3
2
-1 3
4 -2 1
输出
6
\end{verbatim}
二维状态 f[i][j] 表示到第 i 行第 j 个位置的最优值，合法前驱只有 $(i-1,j-1)$ 与 $(i-1,j)$。初值为顶点，非法位置为负无穷。转移
\[f[i][j]=a[i][j]+\max(f[i-1][j-1],f[i-1][j]).\]
上例路径 $2,3,1$ 得6。因为每行只依赖前一行，可压缩成一维。\textbf{从右向左更新}才能确保 f[j-1] 尚未被本行覆盖。代码设置 f[1]=0 作为顶点前的虚拟起点，其他位置为负无穷。时间 $O(n^2)$，额外空间 $O(n)$。
\Code{14_triangle.cpp}
\subsection{每件物品最多选一次（教学例题）}

\practice{P1048}{采药}
时间是容量，药材价值是收益；一维转移按容量逆序。

有 n 件物品，每件体积 c、价值 w，容量为 V，求总体积不超过 V 的最大总价值。$1\le n\le200$，$0\le V\le5000$，$1\le c\le5000$，$0\le w\le10^9$。输入第一行 n、V，后面每行 c、w。
\begin{verbatim}
输入
3 5
2 3
3 4
4 5
输出
7
\end{verbatim}
令 f[i][v] 为只考虑前 i 件物品、容量不超过 v 的最大价值。第 i 件选或不选：
\[f[i][v]=\max(f[i-1][v],\ f[i-1][v-c_i]+w_i),\]
第二项只在 $v\ge c_i$ 时存在。若压缩掉 i 维，容量必须从大到小更新，使 f[v-c] 仍是上一层状态，避免同一件被重复选。复杂度 $O(nV)$，空间 $O(V)$。
\Code{15_knapsack01.cpp}
\subsection{每种物品可以无限选择（教学例题）}

\practice{P1616}{疯狂的采药}
每种可重复选取，按容量正序转移。

输入格式与上一题相同，但现在 n 表示物品种类，每种可以选任意非负整数件。仍要求体积不超过 V、最大化价值，且体积必须为正。采用同样的约束。
\begin{verbatim}
输入
2 7
2 3
3 4
输出
10
\end{verbatim}
选择体积2的两件和体积3的一件，价值10。二维转移的选取项变为 $f[i][v-c_i]+w_i$，它允许当前种类继续使用。因此一维容量循环从小到大，使 f[v-c] 已经包含本种物品。时间 $O(nV)$，空间 $O(V)$。
\Code{16_knapsack_complete.cpp}
\subsection{无限硬币的组合数（教学例题）}

有 n 种\textbf{互不相同的正整数面额}，每种无限个，问恰好凑出 target 的组合数，硬币选择顺序不算不同方案。输出对 $10^9+7$ 取模的结果。$1\le n\le100$，$0\le target\le10^4$，面额不超过 $10^4$。输入第一行 n、target，第二行所有面额。
\begin{verbatim}
输入
3 5
1 2 5
输出
4
\end{verbatim}
四种为 $1+1+1+1+1$、$1+1+1+2$、$1+2+2$、5。令 f[v] 为已处理面额能凑出的组合数。先枚举面额，再正序枚举金额，确保每种组合按固定面额顺序生成一次。f[0]=1 是构造非空方案的起点。时间 $O(n\cdot target)$，空间 $O(target)$。
\Code{17_coin_count.cpp}
\subsection{每组最多选一件（教学例题）}

\practice{P1757}{通天之分组背包}
原题输入包含物品的组号，应先归组，再按整组转移。

有 g 组物品，每组最多选一件，也可以不选。求容量 V 内的最大价值。$1\le g\le100$，$0\le V\le2000$，每组1到20件，体积为正、价值非负。输入第一行 g、V，随后每组先给物品数 k，再给 k 对体积、价值。
\begin{verbatim}
输入
2 5
2
2 3
3 5
2
2 4
4 7
输出
9
\end{verbatim}
第一组选体积3价值5，第二组选体积2价值4。对每组保留旧数组 f，所有选择都从旧状态转移到 next\_f，next\_f 初始复制 f 表示不选。这种写法直接保证同一组至多贡献一件，比混用循环顺序更容易检查。若总物品数为 M，时间 $O(MV)$，空间 $O(V)$。
\Code{18_group.cpp}
\subsection{按成绩和编号排序（教学例题）}

\textbf{题目。}输入$n$名学生的编号和成绩，成绩降序，相同成绩按编号字典序升序。编号是不含空格的字符串，长度1到30且互不相同；$1\le n\le2\times10^5$，成绩0到100。首行n，随后每行编号和成绩，输出排序后的记录。
\begin{verbatim}
输入
4
b02 90
a01 95
a02 90
a03 90
输出
a01 95
a02 90
a03 90
b02 90
\end{verbatim}
\textbf{解答。}结构体数组保存记录，cmp先判断成绩，再用strcmp判断编号。编号数组容量31，多留一格保存字符串结束符。\texttt{\%30s}限制一次最多读入30个字符。比较次数为$O(n\log n)$，每次字符串比较最多扫描30个字符，记录存储为$O(n)$。
\Code{03_rank.cpp}
\subsection{有序去重（教学例题）}
\textbf{题目。}输入$n$个整数，$1\le n\le2\times10^5$、$|a_i|\le10^9$。第一行输出不同值数量，第二行按升序输出这些值。
\begin{verbatim}
输入
6
3 1 3 2 2 1
输出
3
1 2 3
\end{verbatim}
排序后每个等值段保留一个即可，时间$O(n\log n)$。
\practice{P1059}{明明的随机数}
该原题同样要求去重后升序输出，提交时按原题范围与格式读写。
\Code{27_distinct.cpp}
\subsection{追加和撤销（教学例题）}
\textbf{题目。}首行$q$，$1\le q\le10^5$；随后每条命令为ADD x、UNDO或PRINT，$|x|\le10^9$。ADD末尾追加，UNDO删除末尾一项，空时忽略；PRINT输出当前序列，空时输出EMPTY。保证全部PRINT实际输出的元素数不超过$2\times10^5$。
\begin{verbatim}
输入
7
UNDO
ADD 4
ADD 9
PRINT
UNDO
PRINT
UNDO
输出
4 9
4
\end{verbatim}
每次修改仅涉及末尾，vector适合保存操作结果。时间为$O(q+S)$的均摊量级，S为输出元素总数；没有对S的约束时，反复PRINT的输出本身也可能很大。
\Code{33_undo.cpp}
\subsection{从左向右替换（教学例题）}
\textbf{题目。}输入三行s、p、t，各行长度不超过1000，允许空格；s和t可空，p非空。按原串s从左到右查找p，每次选取最靠左的不重叠出现，将其替换为t。新插入的内容不再次参与匹配，输出结果。例如s为aaaa、p为aa、t为b，结果bb；s为aaa时结果ba。
\begin{verbatim}
输入
one two one
one
X
输出
X two X
\end{verbatim}
用start标记原串尚未处理的位置。找到后先追加未匹配片段，再追加t，随后前进p的长度。始终在原串上find，可避免t包含p时重复替换甚至死循环。答案最长为约$10^6$字符；朴素匹配可用$O(|s||p|+|answer|)$估计，具体string查找实现可能优化，本题规模不依赖这种优化。
\Code{34_replace.cpp}
\subsection{三类括号匹配（教学例题）}
\textbf{题目。}输入一行仅由\texttt{()[]\{\}}组成的字符串，长度0到$10^5$；空串合法。若每个右括号能与最近尚未匹配的同类左括号配对，输出YES，否则NO。
\begin{verbatim}
输入
([]){}
输出
YES
另一组输入
([)]
另一组输出
NO
\end{verbatim}

最近出现的未匹配左括号必须先配对，所以使用栈。扫描右括号时先检查非空，再检查类型；扫描结束还必须检查栈是否为空。每个字符最多入栈出栈各一次，时间和最坏空间均为$O(n)$。只比较左右括号总数会误判顺序或嵌套类型。
\Code{28_brackets.cpp}
\subsection{排队取号（教学例题）}
\textbf{题目。}首行$n$，$1\le n\le10^5$；随后每行PUSH x或POP，编号x为int范围内整数。PUSH入队；POP输出并删除队首，空队列输出EMPTY。
\begin{verbatim}
输入
6
POP
PUSH 5
PUSH 2
POP
POP
POP
输出
EMPTY
5
2
EMPTY
\end{verbatim}
维护进入顺序即可。每条命令$O(1)$，总时间$O(n)$，最多保存n个编号。
\Code{29_queue.cpp}
\subsection{保留最大的k个数（教学例题）}
\textbf{题目。}输入n、k及n个整数，$1\le k\le n\le2\times10^5$，数值在int范围内。重复值按不同出现计数，按降序输出最大的k个数。
\begin{verbatim}
输入
5 3
4 1 9 9 2
输出
9 9 4
\end{verbatim}

维护最多k个元素的小根堆。每读入一个数就入堆，超过k个时删除最小值；留下的恰好是已读前缀中最大的k个。反复弹出得到升序，最后逆转。时间$O(n\log(k+1)+k)$，额外空间$O(k)$。若n很小且一次读完，直接sort也是合理选择。
\Code{32_heap.cpp}
\subsection{动态候选值（教学例题）}
\textbf{题目。}首行q，$1\le q\le2\times10^5$；随后每行ADD x、DEL x或GE x，$|x|\le10^9$。ADD加入一个副本，DEL删除一个副本，不存在则忽略；GE输出当前最小的、不小于x的值，不存在输出NONE。
\begin{verbatim}
输入
8
ADD 5
ADD 5
ADD 9
DEL 5
GE 5
DEL 5
GE 5
GE 10
输出
5
9
NONE
\end{verbatim}
使用multiset保留重复值；DEL采用find加erase迭代器；GE用成员lower\_bound。每条操作至多$O(\log(q+1))$，空间$O(q)$。
\Code{30_neighbors.cpp}
\subsection{单词计数（教学例题）}
\textbf{题目。}输入$n$及n个仅含小写字母的单词，每词长度1到30，$1\le n\le10^5$。按字典序输出每个不同单词及次数。
\begin{verbatim}
输入
5
pear apple pear banana apple
输出
apple 2
banana 1
pear 2
\end{verbatim}
读入时对键计数；map已经按键排列，最后顺序遍历即可，不需要再sort。若有m个不同单词，计数约需$O(30n\log(m+1))$时间，存储约$O(30m)$。
\Code{31_frequency.cpp}
\subsection{B3693：数列前缀和4}

\practice{B3693}{数列前缀和 4}
\textbf{题意。}多组非负整数矩阵，每次求矩形和模$2^{64}$，每组输出所有查询答案的异或和。$n,m\le1000$，累计矩阵项数、查询数各不超过$10^6$。
\textbf{解答。}使用unsigned long long建立二维前缀和。C++无符号加减自然按$2^{64}$回绕，正好实现题目取模；不能改用有符号溢出。每组清零第0行和第0列，正文符号求和公式其余部分相同。$O(nm+q)$。下面给完整程序。
\textbf{自测样例（按原题规则构造）。}
\begin{verbatim}
输入
2
1 2 2
18446744073709551615 1
1 1 1 2
1 1 1 1
1 1 1
0
1 1 1 1
输出
18446744073709551615
0
\end{verbatim}
\Code{68_unsigned_prefix.cpp}


\subsection{P8649：k倍区间}

\practice{P8649}{[蓝桥杯 2017 省 B] k 倍区间}
\textbf{题意。}输入n、k和n项正整数，求和能被k整除的非空连续区间数。$n,k,a_i\le10^5$。
\textbf{解答。}区间和是k倍数当且仅当两端前缀余数相等。顺序扫描，每读到余数r，先加以前的cnt[r]，再增加计数。初始cnt[0]=1代表空前缀；次序反过来会把空区间算进去。$O(n+k)$时间、$O(k)$空间，答案用long long。
\textbf{自测样例（按原题规则构造）。}
\begin{verbatim}
输入
5 3
1 2 3 4 5
输出
7
\end{verbatim}
\Code{64_k_intervals.cpp}


\subsection{P1638：逛画展}

\practice{P1638}{逛画展}
\textbf{题意。}n幅画各有1到m的作者编号，求包含全部作者的最短连续区间；同长时输出左端点最小的一组，保证有解。$n\le10^6,m\le2000$。
\textbf{解答。}频率记录窗口内作者出现次数，种类数达到m就尽量缩小左端。每次删除后若该作者计数变0，窗口才失去一类。仅在长度严格变小时更新答案，扫描顺序会保留最小左端。$O(n+m)$。
\textbf{自测样例（按原题规则构造）。}
\begin{verbatim}
输入
7 3
1 2 1 3 2 1 3
输出
2 4
\end{verbatim}
\Code{65_gallery.cpp}


\subsection{P3143：两个展示柜}

\practice{P3143}{[USACO16OPEN] Diamond Collector S}
\textbf{题意。}n颗钻石放进两个柜，每柜最大、最小尺寸差不超过K；每颗只能使用一次，求总展示数最多。$n\le50000$。
\textbf{解答。}先排序。双指针求每个起点的最长合法段，再算后缀中单柜最多能取多少。第一个柜取$[l,r)$，第二个柜从r以后取，贡献len[l]+best[r]。有最优方案可按尺寸把两柜分开，扩大第一段、删去重叠仍不会降低总数；不能直接相加两个可能重叠的最大窗口。$O(n\log n)$。
\textbf{自测样例（按原题规则构造）。}
\begin{verbatim}
输入
7 2
1 2 3 8 9 10 20
输出
6
\end{verbatim}
\Code{66_diamonds.cpp}


\subsection{P10719：黑白格}

\practice{P10719}{[GESP202406 五级] 黑白格}
\textbf{题意。}输入n、m、k及01网格，求至少包含k个黑格的最小子矩形面积；不存在输出0。$n,m\le100$。
\textbf{解答。}枚举上下边界，累加每列黑格数，再用非负数组滑动窗口找和至少k的最短列区间。新增一行只需$O(m)$，避免重新统计每个矩形。$O(n^2m)$时间、$O(nm)$存储。缩小窗口时不要越过右端点；k为正，零列也需要正确删除。
\textbf{自测样例（按原题规则构造）。}
\begin{verbatim}
输入
2 3 3
101
011
输出
4
\end{verbatim}
\Code{67_black_rectangle.cpp}
